\input{preamble.tex}

\title{\huge{Topology}}
\author{\LARGE{Thobias Høivik}}
\date{Fall 2026?} 

\begin{document}
\maketitle

\newpage
\tableofcontents

\newpage
\section*{Intro}
The following are my notes for Topology which I will 
(hopefully) be taking in the fall of 2026. 
As these are my own personal notes I may skip ceirtain 
topics I already understand, be sloppy at times and 
this may cause trouble for anyone who wishes 
to use these notes for themselves so keep that in 
mind if you want to use these. 

\section{Topological spaces and continuity}
We recall from Real Analysis that for a metric 
space $(X,d)$ we have a notion of distance 
and can start describing open sets 
$U \subseteq X$ as sets where for every $x \in U$ 
there exists $\varepsilon > 0$ such that 
\[
    B(x;\varepsilon) \subseteq U, 
\]
meaning that for any $x$ in $U$ there is 
some open ball around it contained entirely in $U$.

In Hilbert spaces, normed spaces, $\mathbb R^n$ (different 
types and instances of metric spaces), etc., this is the usual idea: 
an open set is one where every point has some wiggle room around it.

Topology starts from the observation that many arguments in 
Analysis don't really use the distance $d(x,y)$. They only 
use collections of open sets. 

So instead of letting open sets be generated by some metric, why 
don't just specify which sets are open? 

Intuitively, a topology on a set $X$ is a choice about which 
sets are "open", whatever that may mean. 
Informally a topology tells us which points are "near" eachother 
without really specifing quantitatively what that may mean. 

So we want the least structure possible to be able to talk 
about things like continuity, connectedness, compactness, 
convergence, etc. 

\begin{defn}[Topological space]
    Let $X$ be a set. A \textbf{topology} 
    on $X$ is a collection $\tau \subseteq 
    \mathcal P(X)$ satisfying: 
    \begin{enumerate}
        \item $\emptyset, X \in \tau$.
        \item If $\{U_i\}_{i\in I}$ is any collection of elements 
            in $\tau$, then 
            \[
                \bigcup_{i\in I}U_i \in \tau. 
            \]
        \item If $U_1, \ldots, U_n \in \tau$, then 
            \[
                U_1 \cap \cdots \cap U_n \in \tau. 
            \]
    \end{enumerate}
    The pair $(X,\tau)$ is called a \textbf{topological space}. 
    The elements of $\tau$ are called \textbf{open sets}.
\end{defn}
We will often refer to a topological space by just 
the set if it is clear what the topology on $X$ is.  

The axioms are not arbitrary, but abstracted from metric spaces. 
In metric spaces we already have these properties even if they 
are not explicitly required. 
For instance in $\mathbb R$, 
\[
    \bigcap_{n=1}^\infty \left(-\frac{1}{n},\frac{1}{n}\right) 
    = \{0\}. 
\]
Each interval is open, but $\{0\}$ is not. 

Intuitively, finite intersections preserve some positive amount of 
wiggle room while infinite intersections can squeeze it all away. 

\subsection*{Examples of topological spaces}

\begin{ex}[$\mathbb R$]
    Let $X = \mathbb R$. The usual topology is the 
    collection of subsets $U \subseteq \mathbb R$ such 
    that for every $x \in U$, there exists $\varepsilon > 0$ with 
    \[
        (x - \varepsilon, x + \varepsilon) \subseteq U.
    \]
    This is the topology we know from Real analysis. 
    Examples of open sets include: 
    \[
        (0,1), \ (-3,7), \ \mathbb R, \ \emptyset, 
    \]
    and also things like 
    \[
        (0,1) \cup (2,5)\cup (10, \infty). 
    \]
    Non-examples: 
    \[
        [0,1], \ \mathbb Q, \ \{0\}. 
    \]
    Note that $\mathbb Q$ can be made into a topological space 
    by giving it the subspace topology from $\mathbb R$. 
    More on that later. 
\end{ex}
\begin{ex}[Metric spaces]
    Every metric space $(X,d)$ gives a topology. 
    Define $\tau_d$ by saying taht $U \subseteq X$ is open 
    if for every $x \in U$ there exists $\varepsilon > 0$ such that 
    \[
        B(x;\varepsilon) \subseteq U.
    \]
    Then $\tau_d$ is a topology on $X$. 
    So every metric spaces is a topological space, but 
    note that the opposite implication does not hold which 
    is one of the reasons topology is more general. 
\end{ex}
\begin{ex}[The discrete topology]
    Let $X$ be any set. Let $\tau := \mathcal P(X)$. 
    Then $\tau$ is a topology on $X$ which we call 
    the \textbf{discrete topology}. 
\end{ex}
\begin{ex}[The indiscrete topology]
    Let $X$ be any set. Define $\tau = \{\emptyset, X\}$. 
    This is called the \textbf{indiscrete topology} or 
    \textbf{trivial topology}.
\end{ex}

\begin{defn}[Closed set]
    A subset $C \subseteq X$ is called \textbf{closed} if its 
    complement is open: 
    \[
        C \ \text{is closed} \Leftrightarrow 
        X \setminus C \ \text{is open}.
    \]
\end{defn}

In a topological space, $\emptyset$ and $X$ are closed. 
Arbitrary intersections of closed sets are closed and 
finite unions of closed sets are closed. 

An important nuance is that closed does not mean not open. 
$\emptyset, X$ are both open and closed. In $\mathbb R$ (with the 
usual topology) we have that
$\emptyset, \mathbb R$ are both open and closed while e.g. 
$[0,1)$ is neither open nor closde. 

\subsection*{Continuity}
In Analysis, a function on metric spaces $f : X \to Y$ 
is continuous at $x\in X$ if small changes in $x$ in 
$f(x) \in Y$. Recall the $\varepsilon$-$\delta$ definition: 
\[
    \forall \varepsilon > 0, \ \exists \delta > 0 \ \text{s.t.} \ 
    d_X(x,y) < \delta \Rightarrow d_Y(f(x), f(y)) < \varepsilon.
\]

Since we don't have access to metrics we must rewrite this. 

\begin{defn}[Continuity]
    Let $(X,\tau_X)$ and $(Y,\tau_Y)$ be topological spaces. 
    A function 
    \[
        f : X \to Y
    \]
    is \textbf{continuous} if for every open set $V \in \tau_Y$, 
    the preimage 
    \[
        f^{-1}(V)
    \]
    is open in $X$. 
\end{defn}

\begin{prop}
    Let $(X,d_X)$ and $(Y,d_Y)$ be metric spaces. 
    A function 
    \[
        f : X \to Y
    \]
    is continuous in the $\varepsilon$-$\delta$ sense if 
    and only if for every open set $V \subseteq Y$, the 
    preimage 
    \[
        f^{-1}(V)
    \]
    is open in $X$.
\end{prop}
\begin{proof}
    First suppose $f$ is continuous in the $\varepsilon$-$\delta$ 
    sense. 
    Let $V \subseteq Y$ be open. We want to show that $f^{-1}(V)$ is 
    open in $X$. Take $x \in f^{-1}(V)$. Then 
    $f(x) \in V$. Since $V$ is open, there exists 
    $\varepsilon > 0$ such that 
    \[
        B(f(x) ; \varepsilon) \subseteq V. 
    \]
    By continuity of $f$ at $x$, there exists $\delta > 0$ such that 
    \[
        d_X(x,y) < \delta \Rightarrow d_Y(f(x),f(y)) < \varepsilon. 
    \]
    Thus, if $y \in B(x;\delta)$, then 
    \[
        f(y) \in B(f(x);\varepsilon) \subseteq V. 
    \]
    Therefore 
    \[
        y \in f^{-1}(V). 
    \]
    So 
    \[
        B(x;\delta) \subseteq f^{-1}(V). 
    \]
    Since every point of $f^{-1}(V)$ has a ball around it 
    contained in $f^{-1}(V)$, the set $f^{-1}(V)$ is open. 
    
    Now suppose that for every open set $V \subseteq Y$, the 
    preimage $f^{-1}(V)$ is open in $X$. 
    Let $x \in X$, and let $\varepsilon > 0$. The ball 
    \[
        B(f(x) ; \delta)
    \]
    is open in $Y$. Therefore 
    \[
        f^{-1}(B(f(x);\delta))
    \]
    is open in $Y$. 
    Since $x \in f^{-1}(B(f(x);\delta))$, there exists 
    $\delta > 0$ such that 
    \[
        B(x;\delta) \subseteq f^{-1}(B(f(x);\delta)). 
    \]
    Thus, if 
    \[
        d_X(x,y) < \delta, 
    \]
    then 
    \[
        y \in f^{-1}(B(f(x);\delta)), 
    \]
    so 
    \[
        f(y) \in B(f(x);\delta). 
    \]
    Equivalently, 
    \[
        d_Y(f(x),f(y)) < \varepsilon. 
    \]
    Hence $f$ is continuous at $x$. Since $x$ was arbitrary, 
    $f$ is continuous.
\end{proof}

Suppose $\tau_1$ and $\tau_2$ are two topologies on 
the same set $X$. 

We say $\tau_1$ is \textbf{finer} than $\tau_2$ if 
\[
    \tau_2 \subseteq \tau_1. 
\]
That means $\tau_1$ has more open sets. 
Equivalently, $\tau_2$ is \textbf{coarser} than 
$\tau_1$. For every set $X$: 
\[
    \{\emptyset, X\} \subseteq \tau \subseteq \mathcal P(X). 
\]
So the indiscrete topology is the coarsest topology possible, 
and the discrete topology is the finest topology possible. 

This affects continuity. If the domain topology is finer, 
continuity is easier. If the target topology is finer, continuity 
is harder. 

\begin{ex}
    Let $X$ have the discrete topology and let $Y$ be any 
    topological space. 
    Then every $f : X \to Y$ is continuous. 

    This is because $f^{-1}(V) \subseteq X$ and 
    every subset of $X$ is open. 
\end{ex}
\begin{ex}
    Let $Y$ have the indiscrete topology. Then 
    every function $f :X \to Y$ is continuous. 
    
    The preimages 
    \[
        f^{-1}(\emptyset) = \emptyset, \ 
        f^{-1}(Y) = x
    \]
    are both open in $X$. 
\end{ex}
\begin{ex}
    Let $X$ be a set with two topologies $\tau_1, \tau_2$. 
    Then $\operatorname{id} : (X, \tau_1) \to (X,\tau_2)$ is 
    continuous exactly when $\tau_2 \subseteq \tau_1$. 
\end{ex}

\subsection*{First glimpse of homeomorphisms}

Topology is about spaces up to continuous deformation. 
The correct notion of "same space" is not equality, but 
\textbf{homeomorphism}. 

\begin{defn}[Homeomorphism] 
    Let $X,Y$ be topological spaces. 

    A function $X \to Y$ is a \textbf{homeomorphism} if 
    \begin{enumerate}
        \item $f$ is bijective, 
        \item $f$ is continuous, 
        \item $f^{-1}$ is continuous.
    \end{enumerate}
    If such an $f$ exists, then $X$ and $Y$ are homeomorphic. 
\end{defn}
Intuitively, homeomorphic spaces have the same topological shape. 
Think about those videos of coffee cups and donuts. This is what 
they meant. 

\begin{ex}
    $(0,1)$ and $\mathbb R$ are homeomorphic. 
    One possible homeomorphism is 
    \[
        f : (0,1) \to \mathbb R, \ 
        x \overset{f}{\mapsto} \tan\left(\pi x - \frac{\pi}{2}\right).
    \]
\end{ex}
So $\mathbb R$ is homeomorphic to $(0,1)$ even though $(0,1)$ is 
"small" (bounded) in comparison to $\mathbb R$.

Lastly we remark that openness doesn't need to have 
anything to do with distance. 
In algebraic geometry and commutative algebra we look at 
the prime spectrum of a ring, 
\[
    \operatorname{Spec}(A)
\]
which we endow with a topology called the \textbf{Zariski topology}. 
Closed sets are defined by ideals: 
\[
    V(I) = \{\mathfrak p \in \operatorname{Spec}(A) : 
    I \subseteq \mathfrak p\}. 
\]
So the closed subsets are controlled by equations and ideals. 

\subsection*{Exercises}

\begin{prob}
    Let $X$ be any set. Prove that 
    the discrete topology $\tau = \mathcal P(X)$ is 
    a topology on $X$. 
\end{prob}
\begin{proof}
    $\tau \subseteq \mathcal P(X)$ is ceirtainly satisfied. 
    Furthermore $\emptyset, X \in \tau = \mathcal P(X)$ is 
    also satisfied. 

    Now let $I$ be some index set. Since 
    $\mathcal P(X)$ contains all subsets of $X$ and we 
    necessarily have 
    \[
        \bigcup_{i \in I} U_i \subseteq X
    \]
    for $U_i \in \mathcal P(X)$, and then 
    we have 
    \[
        \bigcup_{i\in I} U_i \in \mathcal P(X) = \tau. 
    \]
    Thus $\tau$ is closed under arbitrary intersections. 

    Let $U_1,\ldots,U_n \in \tau = \mathcal P(X)$. 
    Once again since $U_i \subseteq X$ for all $1 \leq i \leq n$ we 
    have 
    \[
        \bigcap_{i\in[n]} U_i \subseteq X
    \]
    so 
    \[
        \bigcap_{i\in[n]} U_i \in \tau = \mathcal P(X).
    \]
\end{proof}

\begin{prob}
    Let $X$ be any set. Prove that $\tau = \{\emptyset, X\}$ 
    is a topology on $X$. 
\end{prob}
\begin{proof}
    We have $\emptyset, X \in \tau$ by definition and also 
    $\tau \subseteq \mathcal P(X)$. 

    Now notice that the only non-trivial collection 
    of open sets in $\tau$ is the entireity of 
    $\tau$: 
    \[
        \emptyset, X,
    \]
    and this collection is finite. 
    Then $\emptyset \cup X = X \in \tau$ and 
    $\emptyset \cap X = \emptyset \in \tau$. 

    Thus $\tau$ is closed under arbitrary unions 
    and finite intersection adn hence forms a topology 
    on $X$. 
\end{proof}
\begin{prob}
    Let $X$ be a topological space. Prove that arbitrary intersections 
    of closed sets are closed. 
\end{prob}
\begin{proof}
    Let $I$ be some index set so that 
    $C_i$ is closed for every $i \in I$. 

    Then $X \setminus C_i$ is open for every $i \in I$. 
    Then 
    \[
        \bigcup_{i\in I}X \setminus C_i
        = \bigcup_{i\in I} X \cap C_i^c 
    \]
    is open. 
    Taking the complement we get 
    \[
        \left(\bigcup_{i\in I} X \cap C_i^c\right)^c
        = \bigcap_{i\in I} \emptyset \cup C_i 
        = \bigcap_{i \in I} C_i
    \]
    so 
    \[
        \bigcap_{i\in I} C_i  
    \]
    is closed. 
\end{proof}

\begin{prob}
    Let $\tau_1,\tau_2$ be topologies on $X$. Prove that 
    $\operatorname{id} : (X, \tau_1) \to (X, \tau_2)$ is 
    continuous if and only if $\tau_2 \subseteq \tau_1$. 
\end{prob}
\begin{proof}
    Suppose $\tau_2 \subseteq \tau_1$. 
    Let $U \in \tau_2$. Then $U \in \tau_1$ and 
    \[
        \operatorname{id}^{-1}(U) = U \in \tau_1.
    \]
    Thus the preimage of $U \in \tau_2$ is in $\tau_1$. 
    As $U$ was arbitrary we have that $\operatorname{id}$ 
    is continuous. 

    Suppose $\operatorname{id}$ is continuous. Then for any open set 
    $U \in \tau_2$ we have that 
    \[
        \operatorname{id}^{-1}(U) = U \in \tau_1
    \]
    hence 
    \[
        \tau_2 \subseteq \tau_1.
    \]
\end{proof}

\end{document}
